\section[Hauville]{L’indexation Hauville: adapter le langage RAMEAU au droit}

La particularité de la bibliothèque Cujas vient de l'indexation que celle-ci utilise peu commune et connue des autres bibliothèques. En effet, la bibliothèque utilise l'indexation Hauville du nom de sa création de Frédérique Hauville qui a travaillé à Cujas sur le projet de la première bibliothèque numérique. 

La raison derrière ce choix vient du fait qu'au démarrage la bibliothèque utilise des \enquote{mots-clés libres qui généraient des incohérences}. Il n'est pas toujours envisageables de récourir à l'indexation RAMEAU qu'il considérait comme inadapté au droit. En effet, l'indexation Hauville trouve son origine grâce au référentiel RAMEAU, qui est utilisé uniquement quand cela est pertinent et si ce n'est pas le cas le choix a été fait de le déconstruire afin de lui redonner un nom convenant mieux davantage précis. A l'origine cette indexation est validée par la BnF avec une mise en garde de leur part en avançant que cela ne pourra peut-être pas toujours être la solution lorsque la bibliothèque numérique atteindra un nombre conséquent d'ouvrages numérisés. La BnF avait émis l'hypothèse que cela commencerait sans doute à poser un problème à partir de 1 000 documents. 

L'indexation Hauville commence par une analyse documentaire portant sur le titre, la table des matière, la quatrième de couverture ainsi que l'identification des thèmes de l'ouvrage. Ensuite, une fois les mots-clés récupérés ils sont transcrit en utilisant RAMEAU. Si le terme n'existe pas dans l'indexation alors la bibliothèque \enquote{s'appuie sur une forme normalisée qui est la plus courante comme Wikipédia ou vocabulaire des termes juridiques de Cornu par exemple}. Puis, elle construit la vedette pouvant être composé entre deux ou trois éléments en fonction du domaine traité. L'indexation Hauville comprend la vedette de forme (code, travaux préparatoires...), et le concept juridique. 

L'idée est de pouvoir \enquote{relire la vedette à l'envers en mettant un de ou dans le(s)}. Un exemple type de facette pouvant être trouvé sur CujasNum est celle-ci : Eglise -- Administration. Une autre particularité de cette indexation est possibilité d'avoir \enquote{une indexation à la fois générique et spécifique. Il est possible d'avoir plusieurs vedettes construites ayant des liens de hiérarchie entre elles}, comme par exemple : Droit civil -- Obligations et Droit civil -- Contrats.

Au début de cette méthode d'indexation, à Cujas, il était courant qu'elle soit fait par différente personnes car il s'agissait \enquote{de mots-clés apposés dans les notices Dublin Core}. Le problème vient de la multiplication en charge d'ajouter cette indexation qui ne le faisait pas forcément de la même façon, \enquote{entraînant ainsi des incohérences}. Certains pouvaient utiliser par exemple le trait d'union, d'autre le cadratin créant ainsi une nouvelle facette car ce n'est plus le même élément de séparation qui est utilisé. Pour parer ce problème, il a donc été décidé que \enquote{tout le monde indexerait de la même façon, en proposant de recourir à la liste d'autorité RAMEAU sur BNOpale}. Face aux nombreux problèmes rencontrés par la bibliothèque Cujas, la BNF lui suggère de \enquote{faire une liste fermée du vocabulaire utilisé RAMEAU tant qu'il n'y a pas trop d'ouvrages numérisés. A partir de milliers de documents, cela n'est plus possibile}. 

Dans une bibliothèque numérique, l'indexation est un élément important pour naviguer aisément à l'intérieur. Elle provient des métadonnées qui sont \enquote{garants de fonctionnalités de recherche et de navigation satisfaisantes dans une bibliothèque numérique}. L'interopérabilité des données est importante car sans cela la bibliothèque numérique a des difficultés pour communiquer avec les autres bibliothèques pouvant entraîner son évincement progressif de la communauté malgré tous les services novateurs que celle-ci peut proposer.